\documentclass[pdftex, 12pt, a4paper]{report}

% =============================================================
% 1. PACKAGES & GEOMETRY
% =============================================================
\usepackage{mathptmx} % Times New Roman font
\usepackage[top=3cm, bottom=3cm, left=3.5cm, right=3cm]{geometry}
\usepackage[hidelinks]{hyperref}
\usepackage[pdftex]{graphicx}
\usepackage{booktabs}
\usepackage{setspace}
\usepackage{parskip}
\usepackage{amsmath}
\usepackage{amsfonts}
\usepackage{float}
\usepackage{multirow}
\usepackage{algorithm}
\usepackage{algpseudocode}
\usepackage[shortlabels]{enumitem}

% =============================================================
% 2. ISABELLE CODE HIGHLIGHTING (Added for your project)
% =============================================================
\usepackage{listings}
\usepackage{xcolor}


\lstdefinelanguage{Isabelle}{
    keywords={theory, imports, begin, end, type_synonym, definition, lemma, theorem, proof, assume, show, have, thus, qed, unfolding, apply, by, obtain, where, fix, for},
    keywordstyle=\color{blue}\bfseries,
    ndkeywords={nat, int, set, list, bool, pmf},
    ndkeywordstyle=\color{black}\bfseries,
    identifierstyle=\color{black},
    sensitive=true,
    morecomment=[s]{(*}{*)},
    commentstyle=\color{gray}\itshape,
    stringstyle=\color{red}\ttfamily,
    basicstyle=\ttfamily\small,
    breaklines=true,
    frame=single,
    frameround=tttt,
    backgroundcolor=\color{gray!5},
    captionpos=b,
    numbers=left,
    numberstyle=\tiny\color{gray},
    literate=
        {::}{{::}}2              
        {\\in}{{$\in$}}1         
        {~:}{{$\notin$}}2        
        {~=}{{$\neq$}}2          
        {<=}{{$\le$}}1
        {>=}{{$\ge$}}1
        {=>}{{$\Rightarrow$}}2
        {ALL}{{$\forall$}}1
        {EX}{{$\exists$}}1
        {EX!}{{$\exists!$}}2
        {&}{{$\land$}}1
        {|}{{$\lor$}}1
        {-->}{{$\longrightarrow$}}2
}
% =============================================================
% 3. BIBLIOGRAPHY SETTINGS
% =============================================================
% Explicitly using biber backend to avoid errors
\usepackage[backend=biber, style=ieee, sorting=none]{biblatex}
\addbibresource{bib.bib}

% Custom citation formats to match NTU/IEEE guidelines
\DeclareFieldFormat{issn}{}
\DeclareFieldFormat{url}{\iffieldundef{doi}{\url{#1}}{}}
\DeclareFieldFormat[misc]{title}{\mkbibquote{#1}}
\DeclareFieldFormat[online]{note}{\space{#1}}
\DeclareFieldFormat[online]{url}{\space\url{#1}}

\AtBeginBibliography{\normalsize\urlstyle{same}}

% =============================================================
% 4. FORMATTING & LAYOUT (Titles, TOC, Headers)
% =============================================================
\usepackage{titlesec}
\titleformat{\chapter}{\normalfont\scshape\LARGE\centering}{\thechapter}{1em}{}
\titlespacing*{\chapter}{0pt}{-60pt}{1em}
\titleformat{name=\chapter,numberless}[display]{\normalfont\scshape\LARGE\centering}{}{0pt}{}
\titleformat{\section}{\normalfont\scshape\bfseries\Large}{\thesection}{1em}{}
\titleformat{\subsection}{\normalfont\bfseries\scshape\normalsize}{\thesubsection}{1em}{}

\usepackage{tocloft}
\renewcommand{\cftchapfont}{\normalfont}
\renewcommand{\cftsecfont}{\normalfont}
\renewcommand{\cftloftitlefont}{\hfill\normalfont\scshape\LARGE}
\renewcommand{\cftafterloftitle}{\hfill}
\renewcommand{\cftlottitlefont}{\hfill\normalfont\scshape\LARGE}
\renewcommand{\cftafterlottitle}{\hfill}
\setlength{\cftbeforetoctitleskip}{-20pt}
\setlength{\cftbeforeloftitleskip}{-20pt}
\setlength{\cftaftertoctitleskip}{10pt}
\setlength{\cftbeforechapskip}{0pt}
\renewcommand{\cftdotsep}{1}
\renewcommand{\cftchapdotsep}{\cftdotsep}

\usepackage{fancyhdr}
\pagestyle{fancy}
\fancyhf{}
\fancyfoot[R]{\thepage}
\renewcommand{\headrulewidth}{0pt}
\renewcommand{\footrulewidth}{0pt}

\fancypagestyle{plain}{
  \fancyhf{}
  \fancyfoot[R]{\thepage}
  \renewcommand{\headrulewidth}{0pt}
  \renewcommand{\footrulewidth}{0pt}
}

% Strict formatting settings (Prevent hyphenation & breaks)
\setstretch{1.5}
\sloppy
\tolerance=1
\emergencystretch=\maxdimen
\hyphenpenalty=10000
\hbadness=10000

% =============================================================
% 5. DOCUMENT METADATA
% =============================================================
% Adjusted title and code (Removed brackets [])
\title{\uppercase{Formalizing the Isolation Lemma in Isabelle/HOL}}
\newcommand\fypcode{CCDS25-0096} 
\author{Jiang Yifei}
\date{2026}
\newcommand\degree{Bachelor of Computing in Computer Science}

% =============================================================
% 6. MAIN DOCUMENT CONTENT
% =============================================================
\begin{document}

% --- Front Matter ---
% \input{cover.tex} % Uncomment if cover page is needed
\makeatletter
\begin{titlepage}
\begin{center}

% \uppercase{\textbf{\large{Nanyang Technological University}}}
% \\[6cm]
\begin{figure}[!ht]
    \centering
    \includegraphics{fig/NTU Logo.png}
\end{figure}
\vspace{3cm}

\uppercase{\textbf{\fypcode \\[0.3cm]\@title}}
\\[2cm]
\@author

Supervisor: Prof. Tan Yong Kiam
\\[3.5cm]
COLLEGE OF COMPUTING AND DATA SCIENCE

\vfill

Submitted in Partial Fulfilment of the Requirements\\
for the Degree of \degree\\
at Nanyang Technological University
\\[0.8cm]
% by
% \\[0.8cm]





% College of Computing and Data Science\\
\@date
\end{center}
\end{titlepage}
\makeatother

\pagenumbering{roman}
\setcounter{page}{2}

\addcontentsline{toc}{chapter}{Abstract}
\chapter*{Abstract}

Randomized algorithms are a cornerstone of modern theoretical computer science, yet their proofs often rely on informal probabilistic intuition that obscures critical assumptions about independence and measurability.
This report presents a formal verification of the \textbf{Isolation Lemma} -- a fundamental principle ensuring the uniqueness of minimum-weight structures -- within the \textbf{Isabelle/HOL} interactive theorem prover.

We explicitly construct the underlying product probability space for random weight assignments and formalize the conditions for the existence and uniqueness of minimizers over finite sets.
By rigorously defining the auxiliary value $\alpha(x)$ and proving its independence from individual weights $w(x)$, we formally derive the upper bound on the failure probability.
Furthermore, this formalization serves as the theoretical foundation for verifying a randomized parallel algorithm for the \textbf{Perfect Matching problem}.
The results demonstrate how pseudo-formal probabilistic arguments can be transformed into fully machine-checked proofs, clarifying hidden dependencies in algorithmic reasoning.

\pagebreak
\addcontentsline{toc}{chapter}{Acknowledgements}
\chapter*{Acknowledgments}
% \addcontentsline{toc}{chapter}{Acknowledgments}

I would like to express my gratitude to Assitant Professor Tan Yong Kiam for his guidance and support throughout this project. 

TODO: Acknowledgements

% --- Tables of Content/Figures ---
\renewcommand{\contentsname}{\normalfont\scshape\LARGE Table of Contents}
\renewcommand{\listfigurename}{\normalfont\scshape\LARGE List of Figures}

\pagebreak
\tableofcontents

\pagebreak
\addcontentsline{toc}{chapter}{List of Figures}
\listoffigures
% Note: List of Tables is omitted as this is a formalization project.

\cleardoublepage
\pagenumbering{arabic}
\renewcommand{\figurename}{Fig.}

% --- Chapters ---
\chapter{Introduction}
\label{chp:intro}

Randomization has become one of the central paradigms of modern theoretical computer science. 
Over the past few decades, randomized methods have profoundly influenced the design and analysis of algorithms, enabling solutions that are simpler, faster, and even provably unattainable by deterministic methods.

From parallel computation to combinatorial optimization and approximation algorithms, randomness is often treated as a significant resource rather than a secondary heuristic.

Meanwhile, many core results of randomized algorithm theory are supported by probabilistic lemmas, the proofs of which are often highly abstracted.
These proofs often rely on intuitive statements about probability, independence, and minimization, where detailed assumptions are often treated implicitly, and intermediate steps are frequently omitted.

While these proofs are understandable to humans, they present a significant challenge for formal verification within a rigorous logical framework, precisely because they rely on informal probabilistic intuition.
In particular, probabilistic arguments often compress multiple layers of reasoning into concise informal claims, making hidden dependencies difficult to detect.
This gap between informal reasoning and formal justification becomes especially pronounced in the context of randomized algorithms, where correctness is inherently probabilistic rather than absolute. 

From a proof-methodological perspective, this reflects a transition from informal pen-and-paper arguments to pseudo-formal proofs that appear rigorous but still depend on unstated assumptions, and ultimately to fully formal, machine-checkable proofs.
Randomized algorithm proofs are particularly sensitive to this transition, since properties such as independence, measurability, and uniqueness must be stated explicitly rather than justifies by intuition.
The Isolation Lemma provides a compact yet representative case study for this transition, as it brings together these core ingredients in a single probabilistic argument. 

Interactive theorem proving offers a disciplined way to address these problems.
By encoding mathematical definitions and proofs into a formal logical system and verifying them mechanically through explicit reasoning steps, theorem provers provide a level of precision and reliability that cannot be achieved solely through informal reasoning.

In recent years, there has been significant progress in the formalization of mathematics and deterministic algorithms. 
However, formal proofs of randomized algorithms remain relatively underexplored, largely due to the additional complexity introduced by probabilistic reasoning.

This thesis lies at the intersection of these two research directions. 
Its goal is to formally verify a crucial probabilistic lemma for randomized algorithms — the Isolation Lemma — in the Isabelle/HOL theorem proving system.

The Isolation Lemma is a fundamental tool in many randomized algorithms, notably in randomized algorithms for the Perfect Matching problem, where it is used to guarantee the uniqueness of minimum-weight structures with high probability.
Despite its apparent simplicity, the Isolation Lemma involves several implicit assumptions about minimization, independence, and probability spaces.
Explicitly stating and systematically verifying these assumptions is not only a technical challenge but also a valuable opportunity to clarify the logical structure underlying randomized algorithmic reasoning.

Through a complete formalization of the Isolation Lemma and its application to a randomized Perfect Matching algorithm, this work aims to contribute to the broader effort of bringing probabilistic algorithm theory within the scope of machine-checked formal verification.


\section{Background}
\subsection{Randomized Algorithms in TCS}
Randomized Algorithms play a fundamental role in theoretical computer science.

Unlike deterministic algorithms, randomized algorithms incorporate randomness to break symmetry, simplify theoretical constructions, and achieve improved expected performance.
Such algorithms have been widely used in matching problems, parallel computing, approximation algorithms, and the construction of complex combinatorial structures.

In many classical results, the introduction of randomized weights serves as a general technique for enforcing certain "uniqueness" or "isolation" properties.
For example, in several parallel algorithms and combinatorial optimization problems, algorithmic correctness relies on the fact that, after a random perturbation, a minimum structure becomes unique with high probability.

The Isolation Lemma is a canonical instance of this principle. 
Informally, this lemma states that when random weights are assigned independently to elements of a finite universe, there is a high probability that a unique minimum-weight set exists within a given family of subsets.
This lemma provides a theoretical foundation for numerous randomized algorithms and plays a critical role in algorithmic frameworks for problems such as Perfect Matching.

\subsection{Formalization and Theorem Proving}

In theoretical computer science, many fundamental results rely on intricate mathematical arguments and precise logical reasoning.
As theoretical frameworks continue to develop, the scale and complexity of such proofs have grown substantially, making traditional human-based verification increasingly difficult.

Formalized mathematics and interactive theorem proving (ITP) offer a different approach to reasoning about mathematical theories.
In this setting, definitions, theorems, and proofs are encoded within a formal logical system, and every inference step is checked rigorously by a computer.
This provides a level of precision and reliability that is difficult to achieve through informal reasoning alone.

Compared to traditional pen-and-paper proofs, formalized theorem proving exhibits two distinctive characteristics.
First, all implicit assumptions must be made explicit.
These include, for example, the non-emptiness and finiteness of sets, the totality of functions, and the measurability of probability spaces.
Second, the correctness of a proof does not rely on the intuition of the reader, but is guaranteed by the underlying logical framework of the theorem proving system.

In recent years, formalized methods have been successfully applied across many areas in computer science, including program verification, compiler correctness, and the formalization of foundational mathematical theories.
Within theoretical computer science, formal verification has gradually expanded from classical mathematical results to core theorems in algorithm theory and computational complexity.

However, compared to deterministic settings, the formal verification of randomized algorithms poses additional challenges.
Randomized algorithms involve not only complex combinatorial structures, but also probabilistic concepts such as independence, distributions, and expectations.
In a formal setting, these notions must be expressed through explicit constructions of probability spaces, rather than being assumed implicitly.

Against this background, the formal verification of key probabilistic lemmas in randomized algorithms has become an important research direction in interactive theorem proving.
The Isolation Lemma serves as a representative example of such results, combining combinatorial reasoning with probabilistic arguments, and thus provides an ideal case study for formalization.

\section{Motivation}
Although randomized algorithms have achieved significant success in both theory and practice, their correctness often relies on informal probabilistic and combinatorial reasoning.
In traditional pen-and-paper proofs, conclusions such as “a minimum obviously exists”, “this event occurs with probability at most 1/N”, or “the relevant events are independent” are frequently treated as self-evident and left implicit.

In a formal proof setting, however, such assumptions cannot be taken for granted.
All logical dependencies must be stated explicitly and verified.
For instance, the existence of a minimum depends on the non-emptiness and finiteness of the underlying set; probability statements require measurability; and independence cannot be asserted informally, but must be derived from the explicit construction of an underlying probability space.
From a proof-theoretic perspective, this difficulty reflects a broader transition in proof development.
Many arguments in randomized algorithms originate as informal pen-and-paper proofs, and are later refined into pseudo-formal proofs that closely resemble formal reasoning but still rely on implicit assumptions.
Bridging this gap requires transforming such pseudo-formal arguments into fully formal proofs in which all probabilistic and structural dependencies are made explicit.

As a consequence, randomized algorithms and their associated lemmas constitute a particularly challenging class of objects for interactive theorem proving (ITP). 
On the one hand, they occupy a central position in theoretical computer science; on the other hand, they clearly demonstrate the strength of formal verification in clarifying hidden assumptions and exposing implicit reasoning steps.

The Isolation Lemma is particularly well suited to serve as a case study for this transition.
Its proof simultaneously involves probabilistic modeling, independence assumptions, minimization over finite structures, and the emergence of uniqueness through randomization.
As a result, formalizing the Isolation Lemma requires addressing many of the central challenges that arise when translating informal probabilistic reasoning into a fully formal proof.

This project is motivated by these considerations and aims to reconstruct the complete proof of the Isolation Lemma within a fully formalized and rigorous framework.

\section{Project Objectives}
The primary objective of this project is to formally model and prove Isolation Lemma within the Isabelle/HOL theorem proving system.
Beyond establishing the statement of the lemma and its probability bound; this project aims to demonstrate a systematic transition from informal and pseudo-formal probabilistic arguments to fully formal, machine-checkable proofs.

In particular, this involves making explicit the assumptions and constructions that are typically left implicit in pen-and-paper proofs, including the underlying probability space, independence of random variables, and minimization over finite structures.
Through this process, the project seeks to illustrate how pseudo-formal reasoning formed by translation from pen-and-paper proof can be transformed into a rigorous formal proof.

Specifically, this project aims to:
\begin{itemize}
    \item formalize random weight functions together with their corresponding product probability spaces, thereby making probabilistic assumptions explicit;
    \item provide a complete formal proof of the Isolation Lemma, including an explicit upper bound on the probability of non-uniqueness, and a precise characterization of the conditions under which uniqueness fails;
    \item apply the formalized lemma to the correctness analysis of randomized algorithms, illustrating how the resulting formal proof support higher-level algorithmic reasoning.
\end{itemize}

As a concrete application, this project employs the Isolation Lemma to formally verify a key probabilistic argument in a randomized algorithm for the Perfect Matching problem.
This application serves to demonstrate the practical relevance of the formalization and show how isolation-based arguments can be carried through the full transition from informal reasoning to formal verification.

\section{Scope of Work}
The scope of this project is primarily focused on theoretical formalization and proof construction.
In particular, this report covers:

\begin{itemize}
    \item the mathematical formulation of the Isolation Lemma and its formal expression;
    \item the probabilistic modeling of random weight assignments and their associated properties;
    \item the formalization of critical construction steps within the proof;
    \item a representative application of the lemma to Perfect Matching.
\end{itemize}

This work does not address general derandomization techniques for randomized algorithms, nor does it consider implementation-level concerns or performance optimization of specific algorithms.
Moreover, the internal implementation details of the Isabelle/HOL system and its underlying logical foundations are beyond the scope of this report.
\chapter{Preliminaries}
\label{chp:prelims}


\section{Interactive Theorem Proving}


\section{Isabelle/HOL System}

\section{The Isolation Lemma}
The Isolation Lemma, originally introduced by Mulmuley, Vazirani, and Vazirani~\cite{mulmuley1987matching}, is a probabilistic principle used to ensure that, with high probability, a unique solution exists among a family of sets.
It has found widespread applications in randomized algorithms, particularly in problems involving combinatorial search, matching, and derandomization.

\vspace{0.5cm}
\noindent
\textbf{Informal Statement.}
Let $U$ be a finite set (called the \emph{universe}), and let $\mathcal{F} \subseteq \mathcal{P}(U)$ be a non-empty family of subsets of $U$.
Each element $x \in U$ is assigned a weight $w(x)$ independently and uniformly at random from $\{1, 2, \dots, N\}$ for some integer $N \ge 1$.
The total weight of a set $S \in \mathcal{F}$ is defined as"
\[
    w(S) = \sum_{x \in S} w(x)    
\]
The Isolation Lemma states that, with high probability, there is a unique set in $\mathcal{F}$ that minimizes this weight function.

\vspace{0.5cm}
\noindent
\textbf{Theorem (Isolation Lemma)~\cite{mulmuley1987matching}:}
\emph{
    Let $U$ be a finite set with $|U| = n$, and let $\mathcal{F} \subseteq \mathcal{P}(U)$ be a non-empty family of subsets.
    Suppose that each $x \in U$ is assigned an integer weight $w(x)$ drawn independently and uniformly at random from $\{1, 2, \dots, N\}$.
    Then the probability that there exists a unique set $A \in \mathcal{F}$ that minimizers $w(S)$ is at least $1 - \frac{n}{N}$.
}
\[
    \Pr[\exists! A \in \mathcal{F} \text{ minimizing } w(A)] \ge 1 - \frac{n}{N} 
\]

\vspace{0.5cm}
\noindent
\textbf{Proof Idea.}
The proof relies on analyzing the probability that \emph{two or more} sets achieving the same minimum weight.
The key technique is to consider the structure of conflicting pairs and show that their probability of collision is small.
By summing over all such pairs and bounding their contribution, one obtains the upperbound $\frac{n}{N}$ on the failure probability.

\vspace{0.5cm}
\noindent
\textbf{Proof (Joel Spencer~\cite{spencer1995probabilistic}).}
Let $U = \{e_1, e_2, \dots, e_n\}$ be a finite set of $n$ elements, and let $\mathcal{F} \subseteq \mathcal{P}(U)$ be a non-empty family of subset of $U$.
Suppose each element $x \in U$ is assigned a random weight $w(x)$ uniformly from $\{1, 2, \dots, N\}$.

For any $x \in U$, define:
\[
\alpha(x) = \min_{\substack{S \in \mathcal{F} \\ x \notin S}} w(S)- \min_{\substack{S \in \mathcal{F} \\ x \in S}} w(S \setminus \{x\})
\]
This value $\alpha(x)$ depends only on the weights of elements other than $x$ itself, since $x$ is excluded from all weight sums in the definition.
Because $w(x)$ is independently and uniformly chosen from $\{1, 2, \dots, N\}$, the probability that $w(x) = \alpha(x)$ is at most $1/N$.

Now, suppose there are two distinct sets $A, B \in \mathcal{F}$ such that both minimize the total weight $w(S) = \sum_{x \in S} w(x)$.
Let $x \in A \triangle B$ (i.e., an element in the symmetric difference). Without loss of generality, suppose $x \in A \setminus B$.
Then we have:
\[
\begin{aligned}
\alpha(x) &= \min_{\substack{S \in \mathcal{F} \\ x \notin S}} w(S)- \min_{\substack{S \in \mathcal{F} \\ x \in S}} w(S \setminus \{x\})
          &= w(B) - \left( w(A) - w(x) \right)
          &= w(x)
\end{aligned}
\]
So for this $x$, $w(x) = \alpha(x)$. The chance that this happens for any particular $x$ is at most $1/N$, and since there are at most $n$ such $x$, the total probability that some $x$ satisfies this is at most $n/N$.

Thus, with probability at least $1 - \frac{n}{N}$, no such $x$ exists, and the minimum-weight set in $\mathcal{F}$ is unique.

\vspace{0.5cm}
\noindent
\textbf{Remarks.}
\begin{itemize}
    \item The bound $1 - \frac{n}{N}$ improves as $N$ increases. In particular, for $N > 2n$, the probability of uniqueness exceeds $0.5$.
    \item The lemma is non-constructive: it guarantees uniqueness probabilistically, but not which set is unique.
    \item It is widely used in randomized parallel algorithms and in complexity theory, e.g., in the isolation of perfect matchings.
\end{itemize}

\vspace{0.5cm}
\noindent
\textbf{Commentary.}
While the above proof~\cite{spencer1995probabilistic} is concise and elegant, it omits many low-level details that are often assumed in traditional mathematical exposition.
For example, the definition of $\alpha(x)$ implicitly assumes that the minima involved are well-defined, and that comparisons among set weights are meaningful under all random assignments.
Additionally, the reasoning about probabilities (e.g., the independence and uniformity of weights) relies on intuition rather than formal measure-theoretic grounding.
These gaps, though acceptable in informal proofs, present significant challenges when one attempts to formalize such arguments in a theorem prover like Isabelle/HOL.
Thus, part of the motivation of this project is bridge this gap by giving a fully rigorous, machine-checked account of the Isolation Lemma, including all technical conditions and edge cases often left implicit in classical proofs.
In the next chapter, we will present a formalization of this result in Isabelle/HOL, making all implicit assumptions explicit and verifying each step within a constructive logical framework.

\section{Parallel Algorithm for Perfect Matching}

\chapter{Formalizing the Isolation Lemma}
\label{chp:modeling}

The formalization of any randomized algorithm requires a careful translation of intuitive mathematical concepts into rigorous logical definitions.
In this chapter, we will establish the foundational modeling framework for the Isolation Lemma in Isabelle/HOL.

We proceed by constructing the essential components of the lemma layer by layer.

First, we define the underlying universe of the elements and the weight functions, and explain the rationale of choosing natural numbers for weight representation.

Second, we formalize the concept of the "minimizer" and the condition of its uniqueness within a family of sets.

Finally, we construct the probability space using HOL-Probability Library, ensuring that the weights are distributed uniformly and independently.

These definitions act as the foundation upon with the correctness proof in Chapter \ref{chp:proof} is built.
Special attention is given to the design decisions that facilitate proof automation.
In particular, we carefully handel finiteness assumptions and the construction of measure spaces, as these aspects are crucial for enabling automated reasoning in Isabelle/HOL.

\paragraph{Standing assumptions.}
Throughout this chapter, we fix a universe $U$ and a non-empty family of sets $F \subseteq \mathcal{P}(U)$, and assume that $U$ is finite.
For the randomized construction, we fix an integer $N \ge 1$ and sample weights for all $x \in U$ independently and uniformly from $\{1, \dots, N\}$.

These definitions act as the foundation upon which the correctness proof in Chapter ~\ref{chp:proof} is built.
Special attention is given to design decisions that facilitate proof automation, especially regarding the handling of finite sets and measure spaces.


\section{Modeling Weights and Universe}
\label{sec:modeling_weights}

In the standard mathematical formulation of the Isolation Lemma, we consider a universe of elements U.
A weight function, $w$, assigns a numerical value to each element in $U$. 
Typically, these weights are integers chosen uniformly and randomly from a range $\{1, \dots, N\}$.
The weight of a subset $S \subseteq U$ is then defined as the sum of the weights of its constituent elements:
\begin{equation}
    w(S) = \sum_{x \in S} w(x)
\end{equation}

To formalize this in Isabelle/HOL, we avoid defining an explicit "Universe" type.
Instead, we represent the universe implicitly via type variables $'a$.
More precisely, the universe is modeled as an explicit set $U :: 'a~set$.
The type variable $'a$ serves only as a carrier type, while all semantic statements about weights are restricted to elements $x \in U$.
Although weight functions are total by definition in Isabelle, values outside the universe do not play any role in the proof.
We introduce a type synonym to model the weight assignment.
This approach enhances the readability of the subsequent definitions, and ensures a type safety within the proof, which is necessary for the proof assistant.


\begin{lstlisting}[language=Isabelle, caption={Definition of Weight Type and Weight Function}, label={lst:weight_def}]
type_synonym 'a weight = "'a => nat" 
  (* weight maps an element to a natural number *)

definition wt :: "'a weight => 'a set => nat"  where 
  "wt w S = (SUM x \in S. w x)"  
  (* sum of weight *)
\end{lstlisting}

As shown in Listing \ref{lst:weight_def}, we define the type \texttt{'a weight} as a function mapping an element of an arbitrary type \texttt{'a} to a natural number (\texttt{nat}).
Based on this type, the total weight of a set $S$ is defined using the \texttt{wt} function.
 
The choice of a total function is a technical requirement of Isabelle/HOL, where all functions must be defined on the entire carrier type.
Semantically, all arguments and constructions in the sequel will only depend on the values of $w$ restricted to the universe $U$.
This separation between the carrier type and the explicit universe avoids treading the weight function as a partial function, which is not natively supported in Isabelle/HOL.

\subsection*{Design Decision: Natural Numbers vs. Integers}
A critical modeling decision was the choice of the codomain for the weight function.
While the original lemma is often stated using integers, we selected Isabelle's natural number type (\texttt{nat}).
This decision is motivated by two factors:

\begin{itemize}
    \item \textbf{Non-negativity:} In the context of the Isolation Lemma, weight are strictly positive (from $\{1, \dots, N\}$).
        Using \textbf{nat} naturally enforces this constraint without requiring additional invariants to check for negative values.
    \item \textbf{Inductive Reasoning:} The \texttt{nat} type in Isabelle supports powerful induction heuristics which facilitate the automation of arithmetic proofs.
\end{itemize}

\subsection*{Handling Finiteness}
The definition of \texttt{wt} utilizes Isabelle's \texttt{sum} operator(represented as $\sum$ in the math context).
It is crucial to note that in Isabelle/HOL, the summation operator is only mathematically meaningful when the set $S$ is finite.
If $S$ were infinite, the summation would return 0 by definition in the standard library.
This finiteness condition is usually implicit in informal proofs, but must be enforce explicitly in Isabelle/HOL to avoid unintended semantics.

Consequently, throughout our formalization, we do not rely on an implicit assumption that the universe in finite, which is often overlooked in pen-and-paper proofs.
As it is a mandatory step in formal verification, we explicitly enforce the assumption \texttt{finite S} or \texttt{finite U} locally in all lemmas involving weight calculations.
Accordingly, all lemmas involving weight computations in this development explicitly assumes \texttt{finite U} or \texttt{finite S}, making finiteness a first-class requirement rather than an implicit convention.



\section{Defining Minimizers and Uniqueness}
\label{sec:defining_minimizers}

We define "minimizer" as a central concept in the Isolation Lemma.
Intuitively, a minimizer is a set in the family $\mathcal{F}$ that achieves the minimum total weight.
Formally, for a family of subsets $\mathcal{F}$ and a weight function $w$, a set $A$ is a minimizer if:
\begin{equation}
    A \in \mathcal{F} \quad \text{and} \quad \forall B \in \mathcal{F}, \; w(A) \le w(B)
\end{equation}

We formalize this concept defining the predicate \texttt{minimizer F w A}.
This predicate holds true if and only if the set $A$ belongs to the family $F$ and its weight is less than or equal to the weight of any other set $B$ in $F$

\begin{lstlisting}[language=Isabelle, caption={Formal Definition of a Minimizer and Uniqueness}, label={lst:minimizer}]
definition minimizer :: "'a set set => 'a weight =>'a set => bool" where
  "minimizer F w A =
    (A \in F & (ALL B \in F. wt w A <= wt w B))" 
  (* A is one minimizer of F under weight w *)

definition has_unique_minimizer :: "'a set set => 'a weight => bool" where
  "has_unique_minimizer F w =
    (EX! A \in F. minimizer F w A)"
  (* then A is the unique minimizer of F under weight w *)
\end{lstlisting}

As illustrated in Listing \ref{lst:minimizer}, the definition captures the standard mathematical intuition directly.
It is important to note that this definition is purely predicate-based.
The predicate \texttt{minimizer F w A} does not, by itself, assert the existence or uniqueness of such a set.
These properties must be established separately under additional assumptions on the family $F$ and the weight function $w$.
In the setting of formal proof, such distinctions are neccessary to avoid implicitly assuming the existence or uniqueness of minimizers.

The core assertion of the Isolation Lemma concerns the \textbf{uniqueness} of this minimizer.
We define a predicate \texttt{has\_unique\_minimizer} to capture this property formally.
Here, we utilize Isabelle's unique existential quantifier $\exists!$.
The expression $\exists! A. P(A)$ asserts that there exists \textit{exactly one} element $A$ such that the predicate $P(A)$ is true.

Thus, \texttt{has\_unique\_minimizer F w} states that there is exactly one set $A$ in the family $F$ that satisfies the minimizer condition.
This concise definition serves as the target statement for our main probability bound proof in the subsequent chapters.

\subsection*{A Useful Reformulation of Non-Uniqueness}

In the proof of the Isolation Lemma, it is often more convenient to reason about non-uniqueness by working with explicit witnesses.
Accordingly, we will make use of the standard logical equivalence that the minimizer is not unique if and only if there exist two distinct sets $A \neq B$ in $\mathcal{F}$ that both satisfy the minimizer predicate.
This reformulation allows us to fix such a pair of minimizers and derive concrete structural properties in the subsequent proof.

\section{The Alpha Construction}
\label{sec:alpha_construction}

The crux of the Isolation Lemma relies on a specific value associated with each element $x$, denoted as $\alpha(x)$.
This quantity was introduced by Mulmuley et al.\ to isolate the dependency of the total weight on the individual weight $w(x)$.

Mathematically, $\alpha(x)$ is defined as the difference between two minimum values:
\begin{equation}
    \alpha(x)
    =
    \min_{S \in \mathcal{F} \;:\; x \notin S} w(S)
    \;-\;
    \min_{S \in \mathcal{F} \;:\; x \in S}
    w(S \setminus \{x\})
\end{equation}

Intuitively, this measures the "gap" or threshold for $w(x)$.
The first term represents the minimum weight achievable without using $x$.
The second term represents the minimum "residual" weight required to complete a set that must contain $x$.

\subsection*{Formal Definition}
In Isabelle/HOL, we formalize this using the \texttt{LEAST} operator.
The two minimum values are computed independently, and their difference is taken to obtain $\alpha(x)$.

\begin{lstlisting}[language=Isabelle, caption={Definition of Alpha}, label={lst:alpha_def}]
definition alpha :: "('a => nat) => 'a set set => 'a => nat" where
  "alpha w F x =
     (LEAST n. EX S \in F. x ~: S & wt w S = n)
   - (LEAST n. EX S \in F. x \in S & wt w (S - {x}) = n)"
\end{lstlisting}

\subsection*{Well-Definedness and Intended Use}
The use of the \texttt{LEAST} operator requires the existence of witnesses for the corresponding predicates.
In our development, such existence properties are not assumed implicitly, but are discharged later as explicit proof obligations, under suitable conditions on the family $\mathcal{F}$, in particular the existence of minimizers.

At the modeling stage, $\alpha(x)$ is treated as an abstract quantity defined in terms of minimal weights.
No assumptions are made at this point about which sets realize these minimum values.

In the subsequent proof, we will show that once minimizers are known to exist, the \texttt{LEAST} expressions in the definition of $\alpha(x)$ can be simplified to concrete weight values.
This separation reflects a standard discipline in formal verification, where definitions are kept independent of their required assumptions.

\section{Probability Space Construction}
\label{sec:prob_space}

The final component of our modeling framework is the construction of the probability space.
Mathematically, the Isolation Lemma assumes that the weight $w(x)$ for each element $x \in U$ is chosen independently and uniformly at random from the set $\{1,\dots,N\}$.
We assume $N \ge 1$, ensuring that the set $\{1, \dots, N\}$ is non-empty and the uniform distribution is well-defined.

\subsection*{Modeling Challenge: Implicit vs. Explicit Probability}
In standard pen-and-paper proofs, probability is often treated implicitly. We simply state "Let $w$ be a random weight function" and proceed to calculate probabilities of events.
However, in a formal proof assistant like Isabelle/HOL, there is no ambiguity allowed.
We cannot merely assume randomness; we must explicitly construct the underlying probability distribution (as a PMF object).

This introduces a gap between intuition and formalization.
We must define a concrete probability mass function (PMF) over the function space $U \to \{1 \dots. N\}$.
To bridge this gap, we leverage the \texttt{HOL-Probability} library, which formalizes discrete distributions using measure theory.
This construction makes the probability assumptions of the Isolation Lemma explicit objects of reasoning, rather than informal background conditions.


\subsection*{Formalizing the Product Space}
Since the choice for each element $x \in U$ is independent, the joint probability distribution over the universe is the product of the individual marginal distributions.
We construct this in two steps:

\begin{enumerate}
  \item \textbf{Uniform Component:} A uniform distribution over the finite set $\{1, \dots, N\}$.
  \item \textbf{Product Construction:} A product PMF over the finite domain $U$.
\end{enumerate}

This decomposition mirrors the mathematical assumption of independent sampling, but is made explicit here to support formal reasoning about marginals and conditional probabilities.

For the uniform component, we use \texttt{pmf\_of\_set} to represent the uniform distribution over the finite set $\{1, \dots, N\}$.
To combine these marginal distributions, we apply the \texttt{Pi\_pmf} operator, which constructs a product PMF over the domain $U$.

\begin{lstlisting}[language=Isabelle, caption={Construction of the Weight Probability Space}, label={lst:prob_space}]
definition w_pmf :: "'a set => nat => 'a weight pmf" where
  "w_pmf N U = 
    Pi_pmf U (0::nat) (\<lambda>_. pmf_of_set {1..N})"
  (* The probability mass function for random weights *)
\end{lstlisting}

As shown in Listing \ref{lst:prob_space}, \texttt{w\_pmf N U} defines the probability measure over total functions of type \texttt{'a => nat}, where the random choices are made independently on $U$.
The construction uses the \texttt{Pi\_pmf} operator (mathematically $\Pi$) to combine individual uniform distributions into a single product measure.
This explicitly captures the \textbf{independence} assumption, ensuring that the choice of weight for any $x \in U$ does not influence others.

A technical nuance in this definition is the argument \texttt{(0::nat)} passed to \texttt{Pi\_pmf}. 
Since functions in Isabelle/HOL are total, they must return a value even for inputs outside their logical domain ($x \notin U$). Here, \texttt{0} serves as the default value for such out-of-domain queries.
While this implementation detail is necessary to satisfy Isabelle's type system, it does not affect the mathematical validity of the lemma, as our reasoning is strictly confined to the elements within $U$.

By modeling the randomness as a PMF object, probabilistic events become first-class objects in the logic, allowing standard measure theoretic reasnong, such as union bounds, to be carried out within Isabelle/HOL.
\chapter{Proof Construction}
\label{chp:proof}

In this chapter, we present the formal proof of the Isolation Lemma.
The proof follows the structure imposed by the formal development in Isabelle/HOL and proceeds by progressively reducing the failure of uniqueness to a simple probabilistic event.
Unlike pen-and-paper proofs, several steps that are implicit or informal in standard proofs must be mad explicit in formal settings, including existential witnesses, measuability, independence, type correctness, and totality.

In the following sections, we organize the proof as a sequence of structural reductions.
We begin by establishing the existence of minimizers and by reformulating its non-uniqueness in a constructive form.
Then we analyze the auxiliary quantity $\alpha(x)$, which plays a central role in isolating the dependence on individual weights.
Such analysis leads to a characterization of non-uniqueness as a collision event of the form $w(x) = \alpha(x)$.
Finally, we bound the probability of such events using the product probability space constructed in Chapter~\ref{chp:modeling}.

\section{Existence of Minimizers}
\label{sec:existence}

Before analyzing the uniqueness of minimizers, we must first ensure that the notion of a minimizer is well-defined.
In particular, for a given family of sets $\mathcal{F}$ and a weight function $w$, it is necessary to establish that at least one minimizer exists whenever $\mathcal{F}$ is non-empty.

Although the existence of a minimum is intuitively clear when weights take values in the natural numbers, such arguments cannot be taken for granted in the formal setting.
All existence claims must be derived explicitly from previously stated assumptions.
In a formal setting, this step is not merely a formality.
The following definitions and arguments rely on the this predicate as a non-vacuous object, and constructions such as the \texttt{LEAST} operator require the existence of witnesses to be established explicitly.
Isolating the existence of minimizers at this stage allows later proofs to proceed without repeatedly discharging existence conditions.

Formally, we show that for any non-empty family $\mathcal{F}$, there exists a set $A \in \mathcal{F}$ whose total weight is minimal among all elements of $\mathcal{F}$.
That is, $A$ satisfies the minimizer predicate introduced in Chapter~\ref{chp:modeling}.
This result is captured by the following lemma in Isabelle/HOL.

\medskip

\noindent
\textbf{Lemma (Existence of a Minimizer).}
\emph{If $\mathcal{F} \neq \emptyset$, then there exists a set
$A \in \mathcal{F}$ such that $A$ is a minimizer of $\mathcal{F}$ under
the weight function $w$.}

\medskip

This lemma is formalized as \texttt{exist\_minimizer} in the Isabelle development.
The proof constructs a concrete witness $A \in \mathcal{F}$ whose weight is minimal with respect to all other elements of $\mathcal{F}$, thereby ensuring that the minimizer predicate is non-vacuous.

The existence of such a minimizer allows us to reason about minimizers without introducing additional existence assumptions in subsequent arguments.
It also enables a constructive analysis of non-uniqueness, which will be the focus of the next section.

\section{Reformulating the Non-Uniqueness}
\label{sec:non-unique}

The Isolation Lemma concerns the probability that the minimizer of $\mathcal{F}$ under a random weight function $w$ is unique.
In order to reason about the failure of uniqueness in a formal setting, it is necessary to make this notion explicit and structurally usable.

In pen-and-paper proofs, non-uniqueness is often handeled informally by arguing that "there exist at least two different minimizers".
However, in a proof assistant such as Isabelle/HOL, statements involving negation of uniqueness cannot be used directly.
Instead, non-uniqueness must be reformulated as a positive existential statement that provides concrete witnesses.

From a logical perspective, this reformulation eliminates the negation of a unique existential statement to a explicit existential formulation.
Concretely, the minimizer of $\mathcal{F}$ under $w$ is not unique if and only if there exist two distinct sets $A$ and $B$ in $\mathcal{F}$ such that both satisfy the minimizer predicate.
That is, non-uniqueness can be charaterized by the existence of $A \neq B$ with
\[
A \in \mathcal{F}, \quad B \in \mathcal{F}, \quad
\text{and} \quad
\forall S \in \mathcal{F}, \;
w(A) \le w(S) \;\land\; w(B) \le w(S).
\]
Since both $A$ and $B$ satisfy the minimizer predicate, it follows immediately that they attain the same total weight.

This reformulation plays a crucial role in the formal proof.
By providing explicit witnesses $A$ and $B$, it allows subsequent arguments to fix such minimizers and analyze their structural relationship.
In particular, it enables reasoning about elements in the symmetric difference between $A$ and $B$, which forms the bases of the collision characterization developed later.

In the Isabelle/HOL development, this step is realized by unfolding the definition of the unique existential quantifier and reasoning with explicit minimizer predicates.
Although mathematically straightforward, making this reformulation explicit is essential for enabling later arguments that fix concrete minimizers and reason about their set-theoretic structure.


\section{The Alpha Construction}
\label{sec:alpha}

A central ingredient in the proof of the Isolation Lemma is the auxiliary quantity $\alpha(x)$ associated with each element $x \in U$.
Intuitively, $\alpha(x)$ isolates the dependence of the total weight of a set on the individual weight $w(x)$, allowing us to reason about the effect of modifying a single coordinate.

Formally, $\alpha(x)$ is defined as the difference between two minimal weight values determined by the family $\mathcal{F}$.
The first term represents the minimum total weight among all sets in $\mathcal{F}$ that do not contain $x$.
The second term represents the minimum residual weight required to form a set in $\mathcal{F}$ that contains $x$, excluding the contribution of $x$ itself.
This definition was introduced by Mulmuley et al.\ and plays a crucial role in isolating single-coorrdinate events.

\medskip

In our formalization, $\alpha(x)$ is defined using Isabelle's \texttt{LEAST} operator, which selects the least natural number satisfying a given predicate.
At this stage, $\alpha(x)$ is treated as an abstract quantity defined in terms of minimal weights, without assuming that the minima are realized by any particular sets.

\medskip

In the present chapter, we analyze its properties and show how the abstract \texttt{LEAST}-based definition can be simplifies in subsequent arguments, once the existence of minimizers has been derived.

\subsection{Simplifying the Definition of \texorpdfstring{$\alpha(x)$}{alpha(x)}}

The definition of $\alpha(x)$ involves two instances of the \texttt{LEAST} operator, each ranging over a family of sets in $\mathcal{F}$.
While this formalization is convenient at the modeling stage, it is too abstract to be used directly in the subsequent analysis, since the operator \texttt{LEAST} is not computationally usable in formal proof.

Once the existence of minimizers has been established, the two minimization terms appearing in the definition of $\alpha(x)$ can be simplified to concrete weighth values.
More precisely, it can be shown that the minimal values selected by the \texttt{LEAST} operator are realized by specific minimizers, depending on whether one element $x$ is contained in the minimizing set.

\medskip

First, suppose that $B$ is a minimizer of $\mathcal{F}$ under the weight function $w$, and that $x \notin B$.
In this case, $B$ satisfies the defining predicate of the first \texttt{LEAST} expression in the definition of $\alpha(x)$.
As a result, the minimum total weight among all sets in $\mathcal{F}$ that do not contain $x$ is attained by $B$, and the first \texttt{LEAST} expression can be shown equal to the concrete value $W(B)$.

\medskip

Second, suppose that $A$ is a minimizer of $\mathcal{F}$ and that $x \in A$.
Removing $x$ from $A$ yields a residual set whose total weight equals $w(A) - w(x)$.
By showing the $A$ satisfies the defining predicate of the second \texttt{LEAST} expression, the minimum total weight among all residual weights of sets in $\mathcal{F}$ that contain $x$ can be related to the concrete value $w(A \setminus \{x\})$.

\medskip

These observations are formalized by two auxiliary lemmas in the Isabelle development.
They connect the abstract \texttt{LEAST} expressions appearing in the definition of $\alpha(x)$ with concrete weights of minimizers.
Together, they allow $\alpha(x)$ to be rewritten in terms of ordinary arithmetic expressions whenever suitable minimizers are available.

\section{Independence of \texorpdfstring{$\alpha(x)$}{alpha(x)} from Individual Weights}
\label{sec:alpha_independence}

A crucial property of the auxiliary quantity $\alpha(x)$ is that it does not depend on the value of the individual weight $w(x)$.
Instead, $\alpha(x)$ is determined entirely by the weights assigned to those elements other than $x$.
This independence property forms the bridge between the structural analysis of minimizers and the probabilistic argument that follows.

To make this statement precise, consider two different weight functions $w$ and $w'$ that coincide on all elements of $u \setminus \{x\}$, but may differ in the value assigned to $x$.
Changing the value of $w(x)$ shifts the total weight of every set containing $x$ by the same offset, while leaving the relative ordering of weights determined by the remaining elements unchanged.
As a consequence, the minimal values appearing in the definition of $\alpha(x)$ can be shown to be invariant under such changes.

\medskip

Formally, we prove that if two weight functions agree on $U \setminus \{x\}$, then the corresponding values of $\alpha(x)$ are equal.
That is, $\alpha(x)$ depends only on the restrictions of the weight funcion to the complement of $\{x\}$.
This statement is captured by the following lemma in the Isabelle development.

\medskip

\noindent
\textbf{Lemma (Independence of $\alpha(x)$).}
\emph{Let $w$ and $w'$ be two weight functions such that $w(y) = w'(y)$ for all $y \in U \setminus \{x\}$.
Then $\alpha(x)$ computed with respect to $w$ is equal to $\alpha(x)$ computed with respect to $w'$.}

\medskip

This lemma is formalized as \texttt{alpha\_independent\_of\_coord} in Isabelle/HOL.
Its proof relies on the simplification of the \texttt{LEAST} expressions established in the previous section, together with basic congruence properties of addition and inequality on natural numbers.

The independence of $\alpha(x)$ from $w(x)$ is the key technical fact that enables the probabilistic analysis.
Once the weights of all other elements have been fixed, $\alpha(x)$ becomes a constant with respect to the remaining random choice of $w(x)$.
Consequently, the event $w(x) = \alpha(x)$ can be treated as a single-coordinate event in the product probability space, with a simple and uniform probability bound.

\section{Collision characterization via \texorpdfstring{$\alpha(x)$}{alpha(x)}}
\label{sec:collision}

We now establish the key structural link between the failure of uniqueness and the auxiliary quantity $\alpha(x)$.
Specifically, we show that if the minimizer of $\mathcal{F}$ under a weight function $w$ is not unique, then this non-uniqueness much be witnessed by a collision event of the form $w(x) = \alpha(x)$ for some element $x \in U$.

Suppose that the minimizer is not unique.
By the reformulation of non-uniqueness, this implies the existence of two distinct minimizer $A$ and $B$ in $\mathcal{F}$ under the same weight function $w$.
Since $A \neq B$, their symmetric difference in non-empty.
By symmetry, we may assume without loss of generality that there exists an element $x \in A \setminus B$.

\medskip

Because both $A$ and $B$ are minimizers, they have equal total weight, that is, $w(A) = w(B)$.
Applying the simplifications of the \texttt{LEAST} expressioons established in Section~\ref{sec:alpha}, we can express $\alpha(x)$ entirely in terms of the weight of $A$ and $B$.

Since $x \notin B$, the minimum total weight among all sets in $\mathcal{F}$ that do not contain $x$ is realized by $B$
Similarly, since $x \in A$, the minimum residual weight among all sets in $\mathcal{F}$ that contain $x$ in realized by the residual set $A \setminus \{x\}$.
Combining these observations yields
\[
\alpha(x)
= w(B) - w(A \setminus \{x\})
= w(A) - \bigl(w(A) - w(x)\bigr)
= w(x).
\]

Thus, the failure of uniqueness implies the existence of an element $x \in U$ such that the individual weight $w(x)$ coincides with $\alpha(x)$.

\medskip

This argument if formalized in Isabelle/HOL by the lemma \texttt{two\_minimizers\_alpha\_eq}, which shows that whenever two distinct minimizers exist, a corresponding collision event $w(x) = \alpha(x)$ must occur for some $x \in U$.

As a consequence, the event that the minimizer is not unique is contained in the union of the single-coordinate events $\{\, w(x) = \alpha(x) \mid x \in U \, \}$.
This reduction transforms the global failure of uniqueness into a collection of simple local events, paving the way for the probabilistic analysis inthe next section.

\section{Probabilistic Analysis and Final Bound}
\label{sec:probability}

We will now complete the proof of the Isolation Lemma by bounding the probability that the minimizer of $\mathcal{F}$ is not unique under a random weight assignment.
Throughout this section, weights are sampled according to the product probability space constructed in Chapter~\ref{chp:modeling}.
Unlike pen-and-paper arguments, all probabilistic statements here are interpreted with respect to an explicitly constructed probability mass function.

\medskip

By the collision characterization established in Section~\ref{sec:collision}, the event that the minimizer is not unique can be reduced to a union simple single-coordinate events.
More precisely, we have
\[
\{\, \text{the minimizer is not unique} \,\}
\subseteq
\bigcup_{x \in U} \{\, w(x) = \alpha(x) \,\}.
\]
This reduction transforms a global structural failure into a collection of local equality event, each involving only a single random variable.

\medskip

\paragraph{Single-coordinate probability bound.}
Fix an arbitrary element $x \in U$.
By the independence property proved in Section~\ref{sec:alpha_independence}, the auxiliary quantity $\alpha(x)$ depends only on the weights assigned to elements in $U \setminus \{x\}$.
Consequently, once the weights of all other elements have been fixed, $\alpha(x)$ becomes a constant with respect to the remaining random choice of $w(x)$.

Under the product probability space, the random variable $w(x)$ is distributed uniformly over the finite set $\{1 ,\dots,N\}$ and is independent of all other coordinates.
Consequently, the probability that $w(x)$ attains any particular value is exactly $1/N$.
In particular, we obtain the bound
\[
\Pr\bigl[\, w(x) = \alpha(x) \,\bigr] \le \frac{1}{N}.
\]

\medskip

\paragraph{Formalization considerations.}
Although the above argument is mathematically straightforward, its formalization requires several non-trivial steps.
In particular, the event $\{ w(x) = \alpha(x) \}$ must be shown to be a measurable event in the underlying probability space.
This relies on the measurability of the function $\alpha(x)$, which is itself defined using minimum operators over finite families of sets.

Furthermore, the independence arguments depend on the fact that the product probability space assigns weights to to different elements of $U$ independently.
In the Isabelle/HOL development, this property is derived explicitly from the construction of the product PMF using the \texttt{Pi\_pmf} operator, rather than being assumed implicitly.
These considerations account for a substantial portion of the formal proof effort, even though they are typically omitted in pen-and-paper proofs.

\medskip

\paragraph{Union bound and final estimate.}
Since the universe $U$ is finite, we may apply the union bound to the collection of collision events.
Summing the individual bounds yields
\[
\Pr\bigl[\, \text{the minimizer is not unique} \,\bigr]
\le \sum_{x \in U} \Pr\bigl[\, w(x) = \alpha(x) \,\bigr]
\le \frac{|U|}{N}.
\]
Equivalently, the probability that the minimizer is unique is bounded by
\[
\Pr\bigl[\, \text{the minimizer is unique} \,\bigr]
\ge 1- \frac{|U|}{N}
\]

\medskip

This inequality constitutes the statement of the Isolation Lemma.
In the Isabelle/HOL development, the above reasoning is formalized using standard results from the \texttt{HOL-Probability} library, including measurability lemmas for derived random variables and the union bound for finite collections of events.
The final result is captured by the theorem \texttt{isolation\_lemma\_main}, which completes the formal proof.



\chapter{Perfect Matching Problem}
\label{chp:matching}

\chapter{Conclusion and Future Work}
% =============================================================
\section{Summary of Key Findings}
% =============================================================

% =============================================================
\section{Contributions and Impact}
% =============================================================

% =============================================================
\section{Recommendation for Future Work}
% =============================================================

% --- References ---
\begingroup
% Reset formatting locally for bibliography to look better
\tolerance=2000 
\emergencystretch=3em 
\hyphenpenalty=50 
\hbadness=1000 

\addcontentsline{toc}{chapter}{References}
{\fontfamily{ptm}\selectfont 
\normalsize 
\begin{spacing}{1.5}
\printbibliography[title={References}]
\end{spacing}
}
\endgroup

\end{document}